\documentclass[10pt,a4paper]{amsart}
\usepackage[margin=19mm]{geometry}
\usepackage{amsmath,amssymb,mathtools}
\usepackage[scr=rsfs,cal=euler]{mathalfa}
\usepackage{xfrac}
\usepackage[unicode,hidelinks,bookmarks=false]{hyperref}
\newcommand{\ie}{i.e.\ }
\newcommand{\cf}{cf.\ }
\newcommand{\resp}{respectively\ }
\newcommand{\Subsection}{\subsection}
\providecommand{\nobreakdash}{\nobreak\hbox{-}}
\setlength{\emergencystretch}{2em}
\multlinegap0pt
\usepackage[shortlabels]{enumitem}
\usepackage{xcolor}
\usepackage{mathtools}
\usepackage{bm}
\newtheorem{proposition}{Proposition}
\newcommand{\LaguerreL}[2]{L_{#1}^{(#2)}}
\newcommand{\LaguerrecalLjnu}[2]{\mathcal{L}_{#1}^{(#2)}}
\newcommand{\N}{\mathbb{N}}
\newcommand{\Tjnuf}[3]{T_{#1}^{(#2)}#3}
\newcommand{\intervalco}[2]{[#1, #2)}
\newcommand{\intervaloo}[2]{(#1, #2)}
\DeclareMathOperator{\supp}{supp}
\title{Optimal constants of smoothing estimates for quantum harmonic oscillators}
\author{Soichiro Suzuki}
\date{Source: arXiv:2606.03076v1 (CC BY 4.0)}
\begin{document}
\maketitle

\noindent\emph{Source and licence.} Optimal constants of smoothing estimates for quantum harmonic oscillators. Version arXiv:2606.03076v1 (CC BY 4.0), distributed by arXiv under the Creative Commons Attribution 4.0 International licence, http://creativecommons.org/licenses/by/4.0/, as recorded in the arXiv licence metadata for this exact version and shown on the abstract page https://arxiv.org/abs/2606.03076v1. Full licence notice: \texttt{licenses/arxiv-2606.03076-CC-BY-4.0.txt}.\par\smallskip

\noindent\textit{Editorial scope.} Counted unit: the article's proposition behavior (source0.tex lines 551--565). The article's complete original proof of it is delivered with it (source0.tex lines 574--612, one \texttt{proof} environment of 39 lines, which the article places away from the statement). No mathematics is written, repaired or re-derived: the statement and the proof are the article's own lines, in the article's own notation, and no retained line is substituted or re-flowed.  Retained uncounted as the source-local context the proof needs: lines 569--571.  Nothing was omitted from any retained span, so the package carries no omission record.  No retained line contains a citation, so the wrapper carries no bibliography and no bibliography entry.  No retained line contains a figure environment, an image-inclusion command or drawing code, so no image is delivered and the package declares no asset dependency.  Editorial formatting only, in addition: an AMS \texttt{amsart} preamble that restates the article's own declarations and macros used by the retained lines, namely the environments \texttt{proposition} on separate counters, together with the standard packages those lines need.  The retained environments print in this document as Proposition 1 in order of appearance, whereas the article prints the same items with the numbering of its own full document.  The complete native source of the article as distributed in the arXiv e-print for this exact version is bundled unchanged as \texttt{sources/201955/source0.tex}.
	\begin{equation} \label{eq:asymptotic behavior for Laguerre polynomial}
		\LaguerreL{j}{\nu}(r) = \frac{j^{\nu/2-1/4} e^{r/2}}{\pi^{1/2} r^{\nu/2+1/4}} \cos{ (2(jr)^{1/2} - (\nu/2+1/4)\pi )} + O(j^{\nu/2-3/4})
	\end{equation}

\begin{proposition} \label{prop:asymptotic behavior}
Let $\nu \in \intervalco{-1/2}{\infty}$ and $f \in L^1(\intervaloo{0}{\infty})$. Then the following hold.
\begin{enumerate}[label=\textup{(\roman*)}]
\item \label{item:asymptotic behavior for compactly supported f} 
If there exists a compact set $K \subset \intervaloo{0}{\infty}$ such that $\supp f \subset K$, then
\begin{equation} \label{eq:asymptotic behavior for compactly supported f} 
	\lim_{j \to \infty} (4j+2\nu+2)^{1/2} \Tjnuf{j}{\nu}{f} = \frac{2}{\pi} \int_{0}^{\infty} f(r) \, dr .
\end{equation}
\item \label{item:asymptotic behavior for non-negative f} 
 If $f$ is non-negative, then 
\begin{equation} \label{eq:asymptotic behavior for non-negative f} 
	\liminf_{j \to \infty} (4j+2\nu+2)^{1/2} \Tjnuf{j}{\nu}{f} \geq \frac{2}{\pi} \int_{0}^{\infty} f(r) \, dr .
\end{equation}
\end{enumerate}
\end{proposition}

\begin{proof}[Proof of Proposition \ref{prop:asymptotic behavior}]
	First, we prove \ref{item:asymptotic behavior for compactly supported f}. 
	Suppose that $f \in L^1(\intervaloo{0}{\infty})$ has a compact support $K$.
Then, combining $\Gamma(j+1)/\Gamma(j+\nu+1) \sim j^{-\nu}$ and the asymptotic formula \eqref{eq:asymptotic behavior for Laguerre polynomial}, we see that
	\begin{equation} \label{eq:asymptotic behavior for Laguerre function}
	(4j+2\nu+2)^{1/2} \LaguerrecalLjnu{j}{\nu}(r)^2 = \frac{2}{\pi} \sin{(4j^{1/2} r -\nu\pi)} + \frac{2}{\pi} + O(j^{-1/2})
\end{equation}
holds uniformly on $K$ as $j \to \infty$. Therefore, we have
\begin{equation}
(4j+2\nu+2)^{1/2} \Tjnuf{j}{\nu}{f} = \frac{2}{\pi} \int_{r \in K} f(r) \sin{(4j^{1/2} r -\nu\pi)} \, dr + \frac{2}{\pi} \int_{r \in K} f(r) \, dr + O(j^{-1/2})
\end{equation}
as $j \to \infty$.
Using the Riemann--Lebesgue lemma, we obtain
\begin{equation}
\lim_{j \to \infty} (4j+2\nu+2)^{1/2} \Tjnuf{j}{\nu}{f} = \frac{2}{\pi} \int_{r \in K} f(r) \, dr ,
\end{equation}
as desired. 

Next, we prove \ref{item:asymptotic behavior for non-negative f}. 
Suppose that $f \in L^1(\intervaloo{0}{\infty})$ is non-negative and
fix $\varepsilon > 0$ arbitrarily. 
Let $K \subset \intervaloo{0}{\infty}$ be a compact set such that
\begin{equation}
\frac{2}{\pi} \int_{r \in K} f(r) \, dr > \frac{2}{\pi} \int_{r \in \intervaloo{0}{\infty}} f(r) \, dr - \varepsilon / 2 ,
\end{equation}
and let $j_0 \in \N$ be such that 
\begin{equation}
(4j+2\nu+2)^{1/2} \int_{r \in K} f(r) \LaguerrecalLjnu{j}{\nu}(r)^2 \, dr > \frac{2}{\pi} \int_{r \in K} f(r) \, dr - \varepsilon / 2
\end{equation}
holds for every $j \in \N_{\geq j_0}$, which exists thanks to \ref{item:asymptotic behavior for compactly supported f}.
Then we have
\begin{align}
(4j+2\nu+2)^{1/2} \int_{r \in \intervaloo{0}{\infty}} f(r) \LaguerrecalLjnu{j}{\nu}(r)^2 \, dr
&\geq (4j+2\nu+2)^{1/2} \int_{r \in K} f(r) \LaguerrecalLjnu{j}{\nu}(r)^2 \, dr \\
&> \frac{2}{\pi} \int_{r \in K} f(r) \, dr - \varepsilon / 2  \\
&> \frac{2}{\pi} \int_{r \in \intervaloo{0}{\infty}} f(r) \, dr - \varepsilon
\end{align}
for every $j \in \N_{\geq j_0}$. This shows \ref{item:asymptotic behavior for non-negative f}.
\end{proof}

\end{document}
